% ============================================================================
% บทที่ 2 — เตรียมข้อมูลก่อนจัดตารางสอบ
% ============================================================================
\mychapter{เตรียมข้อมูลก่อนจัดตารางสอบ}

\noindent บทนี้อธิบายโครงสร้างไฟล์ข้อมูลที่ระบบรองรับ ประเภทบทบาทของไฟล์ และตัวอย่าง
ข้อมูลที่ถูกต้องและผิดพลาด

\section{ประเภทไฟล์ที่รองรับ}

ระบบรองรับไฟล์ตารางคำนวณทั่วไป เช่น \texttt{XLSX}, \texttt{XLS}, \texttt{CSV},
\texttt{TSV} และ \texttt{ODS} รายการรูปแบบที่รองรับทั้งหมดอยู่ในภาคผนวก A

สำหรับไฟล์ \texttt{CSV} และ \texttt{TSV} ควรบันทึกเป็นรหัส UTF-8 เพื่อให้ภาษาไทย
และรหัสวิชาแสดงผลถูกต้อง ระบบจะอ่านค่าจากเซลล์เป็นข้อความเท่านั้น
ไม่มีการคำนวณสูตรหรือแมโครในไฟล์

\section{โครงสร้างตารางข้อมูลรายวิชา}

ไฟล์ข้อมูลรายวิชาใช้หนึ่งแถวต่อหนึ่งตอนเรียน และหนึ่งคอลัมน์ต่อหนึ่งข้อมูล
คอลัมน์ที่จำเป็นต้องมีคือ:

\begin{center}
\texttt{courseCode} \quad \texttt{courseName} \quad \texttt{sectionNumber} \quad
\texttt{studentGroups}
\end{center}

\begin{longtable}{L{4cm}L{10.5cm}}
\toprule
\textbf{คอลัมน์} & \textbf{คำอธิบาย} \\
\midrule
\texttt{courseCode} & รหัสวิชา เก้าหลัก เช่น \course{061000001} \\
\texttt{courseName} & ชื่อวิชา \\
\texttt{sectionNumber} & หมายเลขตอนเรียน \\
\texttt{studentGroups} & กลุ่มนักศึกษา (สามารถมีได้หลายกลุ่ม คั่นด้วยเครื่องหมายจุลภาค) \\
\bottomrule
\end{longtable}

นอกจากคอลัมน์บังคับ ระบบยังอ่านคอลัมน์เสริมบางรายการ เช่น \texttt{enrolled}
\texttt{capacity} \texttt{instructor} และ \texttt{teachingHours}
(รายการทั้งหมดอยู่ในภาคผนวก A) คอลัมน์อื่นที่ระบบไม่รู้จักจะไม่ถูกนำมาใช้

\section{บทบาทของไฟล์ข้อมูล}

เมื่อนำเข้าไฟล์รายวิชา ต้องระบุบทบาทของไฟล์แต่ละไฟล์:

\begin{description}
  \item[\texttt{จัดอัตโนมัติ}] ระบบจะจัดเวลาสอบให้โดยอัตโนมัติ
  \item[\texttt{จัดไปแล้ว}] ระบบจะคงเวลาสอบเดิมจากไฟล์ไว้
\end{description}

ต้องมีไฟล์อย่างน้อยหนึ่งไฟล์เป็น \texttt{จัดอัตโนมัติ} ระบบจึงจะนำเข้าได้

\begin{caution}
บทบาทของไฟล์ไม่ได้หมายความว่าทุกวิชาในไฟล์นั้นจะแก้ไขเวลาได้เสมอไป
วิชาที่มีเวลาจากไฟล์ต้นทางอาจถูกกำหนดให้แก้ไขไม่ได้
แม้จะอยู่ในไฟล์ \texttt{จัดอัตโนมัติ} ก็ตาม (ดูบทที่ 5)
\end{caution}

\section{ตัวอย่างข้อมูลรายวิชา}

ตัวอย่างไฟล์ \file{managed.xlsx} ที่ถูกต้อง:

\begin{longtable}{L{2.9cm}L{3.9cm}L{3.1cm}L{3.9cm}}
\toprule
\textbf{courseCode} & \textbf{courseName} & \textbf{sectionNumber} & \textbf{studentGroups} \\
\midrule
\course{061000001} & การจัดการอุตสาหกรรมขั้นแนะนำ & 1 & DEMO-A \\
\course{061000002} & การจัดการคุณภาพเบื้องต้น & 1 & DEMO-A \\
\course{061000007} & การวางแผนการผลิต & 1 & DEMO-GA, DEMO-GB, DEMO-GC \\
\bottomrule
\end{longtable}

สังเกตว่า \course{061000001} และ \course{061000002} อยู่ในกลุ่ม \texttt{DEMO-A}
เดียวกัน ซึ่งหมายความว่านักศึกษากลุ่มนี้ต้องเข้าสอบทั้งสองวิชา หากจัดเวลาทับกัน
จะเกิดข้อขัดแย้ง

\section{รหัสวิชาเก้าหลักและการเก็บเลขศูนย์นำหน้า}

รหัสวิชาเป็นข้อความเก้าหลัก ต้องเก็บเลขศูนย์นำหน้าไว้เสมอ ตัวอย่างเช่น
\course{061000001} หากบันทึกเป็นตัวเลข ระบบอาจอ่านเป็น \course{61000001}
ซึ่งผิด

\begin{caution}
ระบบไม่เดาตัวเลขที่หายไป หากรหัสวิชาขาดเลขศูนย์นำหน้า ระบบจะถือว่าเป็นรหัสอื่น
และอาจไม่พบวิชานั้นในเงื่อนไขหรือตารางสอบ ควรตรวจสอบให้รหัสเป็นข้อความเก้าหลักเสมอ
\end{caution}

\section{รูปแบบวันที่และเวลา}

รูปแบบวันที่ที่แนะนำคือ \texttt{YYYY-MM-DD} เช่น \texttt{2026-08-17}
และรูปแบบเวลาคือ \texttt{HH:MM} เช่น \texttt{09:00} ถึง \texttt{12:00}
ระบบแสดงผลวันที่เป็นภาษาไทย (เช่น \texttt{17 ส.ค. 2569}) แต่การป้อนข้อมูล
ควรใช้รูปแบบมาตรฐานสากลเสมอ

\begin{note}
รูปแบบ \texttt{YYYY-MM-DD} ใช้กับข้อมูลในไฟล์ ส่วนช่องวันที่ในหน้าเว็บ
(เช่น ช่วงสอบและวันหยุด) เป็นช่องเลือกวันที่ของเบราว์เซอร์
ซึ่งแสดงรูปแบบตามภาษาและเบราว์เซอร์ที่ใช้ เช่น \texttt{21/08/2026}
ให้เลือกวันที่จากปฏิทินที่กดขึ้นมาตามปกติ
\end{note}

\section{ไฟล์เงื่อนไขการสอบ}

ไฟล์เงื่อนไข (rules) เป็นไฟล์ที่ไม่บังคับ ใช้ระบุวิชาที่ไม่มีสอบ สอบเต็มวัน และสอบพร้อมกัน
รูปแบบของไฟล์เป็นตารางที่แถวแรกเป็นหัวตารางชื่อเงื่อนไข
แต่ละช่องใต้หัวคอลัมน์คือรหัสวิชาที่อยู่ภายใต้เงื่อนไขนั้น
หัวคอลัมน์ที่ระบบรู้จักมีดังนี้:

\begin{longtable}{L{4.5cm}L{10cm}}
\toprule
\textbf{หัวคอลัมน์ในไฟล์} & \textbf{ความหมาย} \\
\midrule
\texttt{ไม่มีสอบกลางภาค} & ระบบจะไม่จัดสอบกลางภาคในส่วนกลางให้วิชาที่ระบุ เงื่อนไขนี้ใช้กับรายวิชาที่ต้องให้ระบบจัดเวลาเท่านั้น และไม่ลบเวลาสอบที่กำหนดมาแล้วจากไฟล์ต้นทาง \\
\texttt{ไม่มีสอบปลายภาค} & ระบบจะไม่จัดสอบปลายภาคให้วิชาที่ระบุ \\
\texttt{ไม่มีสอบ} & ระบบจะไม่จัดสอบทั้งสองภาคให้วิชาที่ระบุ \\
\texttt{สอบสองช่วง} & วิชาที่ระบุต้องสอบเต็มวัน (ต้องตรงกับช่วงเต็มวันที่ตั้งค่าไว้) \\
\texttt{สอบตรงกัน} & วิชาที่ระบุจัดสอบพร้อมกัน (เชื่อมโยงเวลา) \\
\bottomrule
\end{longtable}

\needspace{12\baselineskip}
ตัวอย่างไฟล์ \file{rules.csv} (แถวแรกเป็นหัวตาราง ช่องว่างหมายถึงไม่มีรหัสวิชา):

\begin{center}
\small
\begin{tabular}{L{2.45cm}L{2.45cm}L{1.8cm}L{2.15cm}L{2.15cm}L{2.2cm}}
\toprule
\texttt{ไม่มีสอบกลางภาค} & \texttt{ไม่มีสอบปลายภาค} & \texttt{ไม่มีสอบ} & \texttt{สอบสองช่วง} & \texttt{สอบตรงกัน} & \\
\midrule
\course{061000008} & & & & & \\
 & \course{061000009} & & & & \\
 & & & \course{061000005} & & \\
 & & & & \course{061000003} & \course{061000004} \\
\bottomrule
\end{tabular}
\end{center}

\begin{expected}
ช่องใต้หัวคอลัมน์ \texttt{ไม่มีสอบกลางภาค} มี \course{061000008}
หมายถึงวิชานี้ไม่มีสอบกลางภาค ส่วนคอลัมน์ \texttt{สอบตรงกัน}
ใส่รหัสวิชาต่อกันได้หลายช่องไปทางขวา รหัสทุกช่องในแถวนั้น
(\course{061000003} และ \course{061000004}) จะถูกจัดสอบพร้อมกัน
ความสัมพันธ์แบบสอบตรงกันอาจมีได้มากกว่าหนึ่งคู่ในไฟล์เดียว
\end{expected}

\begin{caution}
การเชื่อมโยงแบบสอบตรงกัน (สอบตรงกัน) จะย้ายเวลาของทุกวิชาที่เชื่อมโยงพร้อมกัน
เมื่อแก้ไขเวลาของวิชาใดวิชาหนึ่ง (ดูบทที่ 5)
\end{caution}

\section{ตัวอย่างข้อมูลที่ผิดพลาด}

ตารางต่อไปนี้แสดงข้อผิดพลาดที่พบบ่อยในการเตรียมข้อมูล:

\begin{longtable}{L{6cm}L{8.5cm}}
\toprule
\textbf{ข้อผิดพลาด} & \textbf{ผลที่เกิดขึ้น} \\
\midrule
รหัสวิชาขาดเลขศูนย์นำหน้า & ระบบไม่พบวิชาในเงื่อนไขหรือตารางสอบ \\
คอลัมน์หัวตารางไม่ตรงชื่อ & ระบบไม่รู้จักคอลัมน์และอาจข้ามข้อมูล \\
ไม่มีกลุ่มนักศึกษาหรือรูปแบบไม่ถูกต้อง & ระบบไม่สามารถตรวจสอบการทับซ้อนของนักศึกษาได้ \\
วันที่หรือเวลาเริ่มต้นอยู่หลังวันที่หรือเวลาสิ้นสุด & ระบบปฏิเสธช่วงเวลาที่ไม่ถูกต้อง \\
\bottomrule
\end{longtable}

\seealso{บทที่ 3 สำหรับการนำเข้าและตรวจสอบ และบทที่ 6 สำหรับการแก้ปัญหาข้อมูล}
